\clearpage\begingroup\NoSourceHead\SourceLeaf{266}
\def\Bsec#1#2{#1\mathord{\overset{\mathrm{s}}{,}}#2}
\def\Barc#1#2{#1\mathord{\overset{\prime\prime}{,}}#2}
\begin{center}
{\Large\bfseries [ASTRONOMISCHE BEOBACHTUNGEN\\AUS DER BRAUNSCHWEIGER ZEIT.]}\par\medskip
{\large\bfseries VERÖFFENTLICHUNGEN UND BRIEFWECHSEL.}\par\medskip
\textbf{[I.]}\par\smallskip
\textbf{[LÄNGENBESTIMMUNG VON BRAUNSCHWEIG, HELMSTÄDT UND\\WOLFENBÜTTEL.]}
\end{center}
\noindent\rule{\linewidth}{.35pt}
\begin{center}\small
[Monatliche Correspondenz. Band X. S. 302--307. 1804 October.\\
Aus v. \textsc{Zach}'s Bericht über die Kgl. Preussische trigonometrische Aufnahme von Thüringen u. s. w.]
\end{center}
\noindent\rule{\linewidth}{.35pt}
\begin{center}\textls[100]{\{1. Braunschweig.\}}\end{center}
\noindent\hspace*{\parindent}\{Der geheime Rath Freyherr von \textsc{Ende} und Dr. \textsc{Gauss} in Braunschweig hatten die Gefälligkeit, unsere Brockensignale in Braunschweig, Helmstädt und Wolfenbüttel zu beobachten. Der geh. R. v. E. hatte seinen eigenen \textsc{Arnold}'schen Chronometer, und dem Dr. \textsc{Gauss} überschickte ich einen dergleichen zu diesem Behufe. Ersterer beobachtete in seiner Wohnung in der Steinstrasse, letzterer an der Südseite der Stadt auf dem Garten des Kaufmanns \textsc{Köppe}, etwa 400 Par. Fuss westlich, und 3500 Fuss südlich vom Andreas-Thurm. Des Freyh. v. \textsc{Ende} Wohnung ist $\Bsec15$ östlich von \textsc{Köppe}'s Garten. Da die Witterung nicht die günstigste war, so musste man bisweilen zu einzelnen Höhen seine Zuflucht nehmen, um die Zeitbestimmung zu erhalten, worin auch wahrscheinlich die kleinen Differenzen zu suchen sind, welche sich in der Bestimmung der Länge von Braunschweig finden. Diese Bestimmungen laufen also:
\clearpage\SourceLeaf{267}\HistoricalHead{262}{L}{PRAKTISCHE UND SPHÄRISCHE ASTRONOMIE. VERÖFFENTLICHUNGEN UND BRIEFWECHSEL.}
\begin{center}\small\setlength{\tabcolsep}{3pt}\renewcommand{\arraystretch}{1.15}
\begin{tabular}{c|rrr|rrr|r}
\hline\hline
1803&\multicolumn{3}{c|}{\shortstack{Mittlere Zeit\\auf Seeberg}}&\multicolumn{3}{c|}{\shortstack{Mittlere Zeit\\in Braunschweig}}&\shortstack{Länge in Zeit\\westl. v. Seeberg}\\\hline\hline
9. August&$9^{\mathrm h}$&$11^{\mathrm m}$&$\Bsec31$&$9^{\mathrm h}$&$10^{\mathrm m}$&$\Bsec{16}4$&$\Bsec{46}7$\\
&&21&$2{,}5$&&20&$16{,}3$&$46{,}2$\\
&&31&$2{,}9$&&30&$16{,}3$&$46{,}7$\\
&&41&$3{,}3$&&40&$16{,}2$&$47{,}1$\\
&&51&$2{,}6$&&50&$15{,}7$&$46{,}9$\\
&10&1&$3{,}4$&10&0&$16{,}3$&$47{,}1$\\\hline
\multicolumn{4}{l}{Anzahl d. Sign. 6}&\multicolumn{4}{r}{Mittel\quad$\Bsec{46}{78}$}\\\hline
15. August&$9^{\mathrm h}$&$1^{\mathrm m}$&$\Bsec{26}3$&$9^{\mathrm h}$&$0^{\mathrm m}$&$\Bsec{40}9$&$\Bsec{45}4$\\
&&10&$27{,}1$&&9&$40{,}9$&$46{,}2$\\
&&20&$26{,}5$&&19&$41{,}0$&$45{,}5$\\
&&30&$26{,}1$&&29&$41{,}1$&$45{,}0$\\
&&40&$26{,}4$&&39&$41{,}5$&$44{,}9$\\
&&50&$26{,}7$&&49&$41{,}8$&$44{,}9$\\\hline
\multicolumn{4}{l}{Anzahl d. Sign. 6}&\multicolumn{4}{r}{Mittel\quad$\Bsec{45}{32}$}\\\hline
\end{tabular}
\end{center}
\begin{center}\small\setlength{\tabcolsep}{3pt}\renewcommand{\arraystretch}{1.15}
\begin{tabular}{c|rrr|rrr|r}
17. August&$9^{\mathrm h}$&$0^{\mathrm m}$&$\Bsec{24}0$&$8^{\mathrm h}$&$59^{\mathrm m}$&$\Bsec{36}3$&$\Bsec{47}7$\\
&&10&$24{,}4$&9&9&$36{,}1$&$48{,}3$\\
&&20&$24{,}5$&&19&$36{,}3$&$48{,}2$\\
&&30&$24{,}3$&&29&$35{,}8$&$48{,}5$\\
&&40&$24{,}5$&&39&$36{,}0$&$48{,}5$\\
&&50&$25{,}0$&&49&$36{,}4$&$48{,}6$\\
&10&0&$25{,}4$&&59&$36{,}5$&$48{,}9$\\\hline
\multicolumn{4}{l}{Anzahl d. Sign. 7}&\multicolumn{4}{r}{Mittel\quad$\Bsec{48}{39}$}\\\hline
18. August&$6^{\mathrm h}$&$20^{\mathrm m}$&$\Bsec{18}0$&$6^{\mathrm h}$&$19^{\mathrm m}$&$\Bsec{28}3$&$\Bsec{49}7$\\
&&30&$18{,}4$&&29&$28{,}3$&$50{,}1$\\
&&40&$18{,}9$&&39&$28{,}6$&$50{,}3$\\
&&50&$18{,}6$&&49&$28{,}3$&$50{,}3$\\
&7&0&$18{,}6$&&59&$28{,}3$&$50{,}3$\\\hline\hline
\multicolumn{4}{l|}{Anzahl d. Sign. 5}&\multicolumn{3}{l}{Am 18. August}&$\Bsec{50}{14}$\\
\multicolumn{4}{l|}{\phantom{Anzahl}»\quad»\quad»\quad7}&\multicolumn{3}{l}{\phantom{Am}»\quad17.\quad»}&$48{,}39$\\
\multicolumn{4}{l|}{\phantom{Anzahl}»\quad»\quad»\quad6}&\multicolumn{3}{l}{\phantom{Am}»\quad15.\quad»}&$45{,}32$\\
\multicolumn{4}{l|}{\phantom{Anzahl}»\quad»\quad»\quad6}&\multicolumn{3}{l}{\phantom{Am}»\quad9.\quad»}&$46{,}78$\\\hline
\multicolumn{4}{l}{Anzahl d. Sign. 24}&\multicolumn{4}{r}{Mittel\quad$\Bsec{47}{66}$}
\end{tabular}
\end{center}
\begingroup\renewcommand{\thefootnote}{}\footnotetext{[Handschriftliche Bemerkung von \textsc{Gauss} in seinem Exemplar:]
\[
\begin{array}{lrrr}
\text{\textsc{Köppens} Garten}&9.&5&\Bsec{50}{30}\\
&10.&7&47{,}70\\
&15.&7&51{,}21\\
&17.&7&46{,}07\\
&18.&8&48{,}11\\
&22.&5&51{,}30\\
&25.&7&49{,}30\\\cline{2-4}
&&46&\Barc{49}{14}.
\end{array}
\]
}\endgroup
\clearpage\SourceLeaf{268}\HistoricalHead{263}{R}{LÄNGENBESTIMMUNG VON BRAUNSCHWEIG, HELMSTÄDT UND WOLFENBÜTTEL.}
\noindent\hspace*{\parindent}Nach diesen Beobachtungen liegt die Wohnung des Freyh. von \textsc{Ende} in Braunschweig $32^{\mathrm m}\,\Bsec{47}{34}$ von Paris, oder geographische Länge von Ferro $28^0\,11'\,50''$, und Kaufmann \textsc{Köppe}'s Garten $32^{\mathrm m}\,\Bsec{48}{84}$\textsuperscript{[*)]} oder von Ferro $28^0\,12'\,\Barc{12}5$.
\begingroup\renewcommand{\thefootnote}{[*)}\footnotetext{Von \textsc{Gauss}' Hand in seinem Exemplar verbessert in] $32^{\mathrm m}\,45^{\mathrm s}\,86$.}\endgroup

Auch hier ergiebt sich dieselbe Probe wie bey Magdeburg. Dr. \textsc{Gauss} beobachtete nämlich die Bedeckung desselben Sterns $\varepsilon$ im Widder vom Monde\textsuperscript{[**)]}. Stellen wir die vom Can. \textsc{David} berechneten wahren Conjunctions-Zeiten abermahls zusammen, und leiten die Längenunterschiede daraus ab, so kommt aus dieser Beobachtung der Längenunterschied von Braunschweig mit Paris aus der Beobachtung
\begingroup\renewcommand{\thefootnote}{[**)}\footnotetext{Im Briefe an \textsc{Olbers} vom 12. August 1803 schreibt \textsc{Gauss}:] Am $9^{\mathrm{ten}}$ habe ich auch Austritt von $\varepsilon$ Arietis beobachtet $10^{\mathrm h}\,38^{\mathrm m}\,\Bsec{39}5$ m[ittlere] Z[eit], wie ich glaube \textls[120]{sehr} gut; von \textsc{Ende} sah ihn mit einem nur sehr nothdürftig aufgestellten oder vielmehr aufgehängten, vielleicht auch an sich schwächern Fernrohr $8^{\mathrm s}$ später. Sollten Sie den Austritt auch beobachtet [haben], so bitte ich um die Communication. [Vgl. auch Monatl. Corresp., Band VIII, S. 468.]}\endgroup
\begin{center}
\begin{tabular}{lr}
von Magdeburg&$32^{\mathrm m}\,\Bsec{45}5$\\
\phantom{von }Wien&$32\quad50{,}0$\\
\phantom{von }Prag&$32\quad40{,}5$\\
\phantom{von }Leipzig&$32\quad39{,}3$\\
\phantom{von }Danzig&$32\quad50{,}0$\\\cline{1-2}
\phantom{von }Mittel&$32^{\mathrm m}\,\Bsec{45}1$
\end{tabular}
\end{center}
\noindent welche astronomische Bestimmung $\Bsec2{99}$\textsuperscript{[***)]} von der terrestrischen abweicht. Merkwürdig ist, dass die Braunschweiger und Magdeburger Beobachtung dieser Sternbedeckung den Längenunterschied für beyde Städte am genauesten gibt. Noch eine andere Prüfung dieser Länge machte Dr. \textsc{Gauss}, indem er auf \textsc{Köppe}'s Garten Azimuthe des Brocken mit seinem Sextanten gemessen hatte. Verbunden mit der Polhöhe $52^0\,15'\,30''$ berechnete er daraus den Längenunterschied $\Bsec{22}9$. Da Seeberg $\Bsec{27}0$ östlich von dem Brocken liegt, so folgt daraus $\Bsec{49}9$ für Seeberg und Braunschweig. Die terrestrische Bestimmung weicht hiervon $\Bsec23$ ab.
\begingroup\renewcommand{\thefootnote}{[***)}\footnotetext{Von \textsc{Gauss}' Hand in seinem Exemplar verbessert in] $\Bsec07$.}\endgroup

Die Breite seiner Wohnung bestimmte der Freyherr v. E. aus vielen Circum-Meridianhöhen der Sonne folgendermassen:
\clearpage\SourceLeaf{269}\HistoricalHead{264}{L}{PRAKTISCHE UND SPHÄRISCHE ASTRONOMIE. VERÖFFENTLICHUNGEN UND BRIEFWECHSEL.}
\begin{center}
\begin{tabular}{rrrr}
1803&6.&Aug.&$52^0\,15'\,\Barc{47}6$\\
»&7.&»&$55{,}1$\\
»&16.&»&$55{,}2$\\\hline
\multicolumn{3}{l}{Mittel}&$52^0\,15'\,\Barc{52}6$.
\end{tabular}
\end{center}
\noindent\hspace*{\parindent}Dr. \textsc{Gauss} hat die Polhöhe auf \textsc{Köppe}'s Garten im vorigen Sommer $52^0\,15'\,35''$ gefunden; fast dasselbe folgt aus der Reduction der Polhöhe seiner vorigen Wohnung, die er $52^0\,16'\,5''$ bestimmt hatte und die $30''$ nördlicher liegt. Ich habe im Sept[ember] 1800 diese Breite im Hotel d'Angleterre $52^0\,15'\,43''$ beobachtet (M. C., II. Bd., S. 562). Daraus würde die Breite vom Dr. \textsc{Gauss} etwas kleiner werden, hingegen nach obiger Bestimmung des Freyh. v. \textsc{Ende} in seiner Wohnung etwas grösser.\}
\begin{center}\textls[100]{\{2. Helmstädt.\}}\end{center}
\noindent\hspace*{\parindent}\{Den 18. August [1803] verfügten sich der Geh. R. von \textsc{Ende} und Dr. \textsc{Gauss} nach Helmstädt, und, nachdem sie daselbst in des Hofraths \textsc{Pfaff} Garten ein Dutzend correspondirende Sonnenhöhen genommen hatten, beobachteten sie hierauf des Abends am 19. Aug. folgende Brocken-Signale:
\begin{center}\small\setlength{\tabcolsep}{3pt}\renewcommand{\arraystretch}{1.15}
\begin{tabular}{c|rrr|rrr|r}
\hline\hline
1803&\multicolumn{3}{c|}{\shortstack{Mittlere Zeit\\auf Seeberg}}&\multicolumn{3}{c|}{\shortstack{Mittlere Zeit\\in Helmstädt}}&\shortstack{Länge in Zeit\\östl. v. Seeberg}\\\hline\hline
19. August&$6^{\mathrm h}$&$0^{\mathrm m}$&$\Bsec{16}5$&$6^{\mathrm h}$&$1^{\mathrm m}$&$\Bsec{24}7$&$1^{\mathrm m}\quad\Bsec82$\\
&&20&$16{,}2$&&21&$24{,}9$&$8{,}7$\\
&&30&$16{,}2$&&31&$24{,}6$&$8{,}4$\\
&&40&$16{,}0$&&41&$25{,}1$&$9{,}1$\\
&&50&$16{,}5$&&51&$25{,}1$&$8{,}6$\\
&7&0&$16{,}7$&7&1&$25{,}2$&$8{,}5$\\
&9&10&$16{,}6$&9&11&$26{,}9$&$10{,}3$\\
&&20&$16{,}2$&&21&$26{,}8$&$10{,}6$\\
&&30&$16{,}6$&&31&$26{,}7$&$10{,}1$\\
&&40&$16{,}7$&&41&$26{,}9$&$10{,}2$\\
&&50&$17{,}1$&&51&$26{,}7$&$9{,}6$\\\hline
\multicolumn{4}{l}{Anzahl d. Sign. 11}&\multicolumn{4}{r}{Mittel\quad$1^{\mathrm m}\quad\Bsec9{30}$}
\end{tabular}
\end{center}
\clearpage\SourceLeaf{270}\HistoricalHead{265}{R}{LÄNGENBESTIMMUNG VON BRAUNSCHWEIG, HELMSTÄDT UND WOLFENBÜTTEL.}
\noindent\hspace*{\parindent}Den 21. August wurden correspondirende Höhen im Gasthofe zum Erbprinzen, und hierauf folgende Signale beobachtet:
\begin{center}\small\setlength{\tabcolsep}{3pt}\renewcommand{\arraystretch}{1.15}
\begin{tabular}{c|rrr|rrr|r}
\hline\hline
1803&\multicolumn{3}{c|}{\shortstack{Mittlere Zeit\\auf Seeberg}}&\multicolumn{3}{c|}{\shortstack{Mittlere Zeit\\in Helmstädt}}&\shortstack{Länge in Zeit\\östl. v. Seeberg}\\\hline\hline
21. August&$9^{\mathrm h}$&$10^{\mathrm m}$&$\Bsec56$&$9^{\mathrm h}$&$11^{\mathrm m}$&$\Bsec{14}5$&$1^{\mathrm m}\quad\Bsec89$\\
&&20&$5{,}5$&&21&$14{,}2$&$8{,}7$\\
&&30&$5{,}6$&&31&$14{,}6$&$9{,}0$\\
&&40&$5{,}8$&&41&$14{,}3$&$8{,}5$\\
&&50&$6{,}0$&&51&$14{,}4$&$8{,}4$\\
&10&0&$5{,}5$&10&1&$14{,}5$&$9{,}0$\\\hline
\multicolumn{4}{l}{Anzahl d. Sign. 6}&\multicolumn{4}{r}{Mittel\quad$1^{\mathrm m}\Bsec8{75}$}
\end{tabular}
\end{center}
\noindent\hspace*{\parindent}Da des Hofraths \textsc{Pfaff} Garten, und der Gasthof zum Erbprinzen ungefähr in demselben Meridian liegen, so kann man für den Längenunterschied zwischen Seeberg und Helmstädt setzen $1^{\mathrm m}\Bsec9{03}$, welches von Paris macht $34^{\mathrm m}\Bsec{44}{03}$, oder Länge von Ferro $28^0\,41'\,\Barc0{45}$.

Den 19. August wurden in des Hofraths \textsc{Pfaff} Garten acht Circum-Meridianhöhen der Sonne beobachtet, welche für die Polhöhe dieses Ortes gaben $52^0\,13'\,\Barc{37}7$. Acht dergleichen Höhen den 21. Aug. im Gasthofe zum Erbprinzen genommen, gaben die Polhöhe $52^0\,13'\,\Barc{51}5$. Auf meiner Harzreise im Jahr 1793 fand ich durch einige Winkel-Beobachtungen auf dem Brocken die Breite durch Interpolation für die Hauptkirche in Helmstädt $52^0\,12'\,58''$, für die Länge $28^0\,40'\,10''$, welches für die Methode, nach welcher ich diese Bestimmungen gemacht und berechnet hatte (Berl. Astr. Jahrb., 1799, S. 141) immer sehr genau ist.\}
\begin{center}\textls[100]{\{3. Wolfenbüttel.\}}\end{center}
\noindent\hspace*{\parindent}\{In Wolfenbüttel wurden den 25. August in des Drosten von \textsc{Rodenberg} Pavillon, nahe beim Schlosse, sechzehn Paar correspondirende Sonnenhöhen, und hierauf folgende Feuer-Signale beobachtet:
\HistoricalFoot{XI\textsubscript{1}.}{34}
\clearpage\SourceLeaf{271}\HistoricalHead{266}{L}{PRAKTISCHE UND SPHÄRISCHE ASTRONOMIE. VERÖFFENTLICHUNGEN UND BRIEFWECHSEL.}
\begin{center}\small\setlength{\tabcolsep}{3pt}\renewcommand{\arraystretch}{1.15}
\begin{tabular}{c|rrr|rrr|r}
\hline\hline
1803&\multicolumn{3}{c|}{\shortstack{Mittlere Zeit\\auf Seeberg}}&\multicolumn{3}{c|}{\shortstack{Mittlere Zeit\\in Wolfenbüttel}}&\shortstack{Länge in Zeit\\westl. v. Seeberg}\\\hline\hline
25. August&$8^{\mathrm h}$&$59^{\mathrm m}$&$\Bsec{41}7$&$8^{\mathrm h}$&$58^{\mathrm m}$&$\Bsec{54}8$&$\Bsec{46}9$\\
&9&9&$41{,}6$&9&8&$53{,}9$&$47{,}6$\\
&&19&$41{,}5$&&18&$54{,}0$&$47{,}5$\\
&&29&$41{,}8$&&28&$54{,}1$&$47{,}7$\\
&&39&$41{,}9$&&38&$54{,}2$&$47{,}7$\\
&&49&$41{,}9$&&48&$54{,}6$&$47{,}3$\\
&&59&$42{,}1$&&58&$54{,}3$&$47{,}8$\\\hline
\multicolumn{4}{l}{Anzahl d. Sign. 7}&\multicolumn{4}{r}{Mittel\quad$\Bsec{47}{50}$}
\end{tabular}
\end{center}
\noindent welches von Paris östliche Länge $32^{\mathrm m}\Bsec{47}5$ und die Länge von Ferro gibt $28^0\,11'\,\Barc{52}5$. Zehn Meridianhöhen der Sonne in Wolfenbüttel gaben die Polhöhe $52^0\,9'\,29''$. Meine Bestimmung von Wolfenbüttel vom Brocken aus und durch genommene Winkel interpolirt, gibt Breite $52^0\,8'\,44''$, Länge $28^0\,11'\,39''$.\}
\clearpage\endgroup
